% The only source of this talk. `make TALK={{SLUG}}` builds slides, notes,
% script, handout, transparencies and article from it. See README.md.
%   - Text outside frames appears only in the article.
%   - Every frame carries a \note{...} for the presenter.
%   - Numbers on slides come from generated files (data/numbers.tex), not typing.
\documentclass[
  beameroptions={aspectratio=169,ignorenonframetext,11pt},
  articleoptions={11pt,a4paper},
]{beamerswitch}
\usepackage{yjtalk}

\title{{{TITLE}}}
\subtitle{One sentence the audience should repeat afterwards}
\author{{{AUTHOR}}}
\institute{Venue or event}
\date{{{DATE}}}

\begin{document}

\begin{frame}[plain]
  \titlepage
  \note{Opening line. Say who this is for and what they will be able to do
    at the end. State the total time.}
\end{frame}

\mode<article>{\maketitle}

\section{Question}

Article-only prose goes here, between frames. It carries the argument in full
sentences for readers who were not in the room.

\begin{frame}{State the finding as a full sentence}
  \begin{itemize}
    \item One message per frame; the title is the message.
    \item<2-> Reveal steps with overlays; the handout and script collapse them.
  \end{itemize}
  \note{What to stress, what to point at, and the transition: name the
    question the next frame answers.}
\end{frame}

\begin{frame}{What to remember}
  \begin{enumerate}
    \item First takeaway.
    \item<2-> Second takeaway.
  \end{enumerate}
  \note{Closing line, then stop talking and take questions.}
\end{frame}

\end{document}
